\cotefile{156-9}
\scanpage{26}
\mspage{22}
% archive sheet number: 25 (upper left)

est \unsure{indispensable}, si on veut pouvoir \unsure{composer|comparer} tous les \uline{\unsure{intervalles}} d'une figure complète $\unsure{\Phi_T|\varphi_T}$ $(T\in\operatorname{Prop}(J))$ telles que
\[
\{J_{<a},\ J_{>a}\}\qquad\text{si }a\in J
\]
ou
\[
J_{<a},\ J_{a,b},\ J_{>b}\qquad\text{si }a<b,\ a,b\in J.
\]

\deleted{Dans} Le cas où $I=\{\omega_0,\omega_1\}$ correspond au cas où $J=\varnothing$, qu'il ne faut nullement exclure ! On trouve alors $\unsure{M|\mathcal M}$ réduit au seul élément $[J]$. On trouve la maquette réduite à une seule \unsure{str.t|Ar.t|Art}, mais visualisée comme figurant un \uline{intervalle ouvert indivisible} : l'intervalle ouvert indivisible est représenté par l'ens ordonné vide !

Si on veut une \unsure{formulation générale} commune, il faut dire qu'on se donne
\[
(I,\Omega_0;\Omega_1)
\]
avec
\[
I,\ \text{ens ordonné},\qquad\Omega_0,\Omega_1\subset I
\]
et les conditions
\margin{donc $\operatorname{Card}\Omega_0,\operatorname{Card}\Omega_1,\leq1$ et il y a quatre cas à distinguer}
\[
\left\{
\begin{aligned}
\text{a)}\quad&\deleted{\unsure{x}}\ x\in\Omega_0\deleted{\unsure{x}}\Longrightarrow x\text{ est un plus petit él. de }I\\
\text{a')}\quad&y\in\Omega_1\Longrightarrow y\ \text{\rule{2.7cm}{0.35pt}\ grand \rule{2.5cm}{0.35pt}}\\
\text{b)}\quad&x\in\Omega_0,\ y\in\Omega_1\Longrightarrow x\neq y.
\end{aligned}
\right.
\]

\scanpage{27}
\mspage{23}
% archive sheet number: 26 (upper left)

Tant qu'il ne s'agit que de décrire un magasin, et y faire une théorie des préfigures, des parties constructibles et des supports, \inserted{et cosupports,} \unsure{qui en dérive}, il n'y a pas à faire d'autres hypothèses sur $J$.

Quand on regarde les préfigures \unsure{se subordonnant|se dessinant} à un drapeau, on \deleted{peut prendre} a intérêt (pour l'énoncé des résultats principaux) à prendre un drapeau qui contienne $\partial J\overset{\mathrm{df}}=\Omega_0\cup\Omega_1$.

Il y a lieu de regarder la structure des \inserted{sous}magasins constructibles induits sur des parties constructibles \inserted{(\unsure{par ex. des supports})} elles-mêmes à des préfigures constructibles \inserted{$\unsure{\Phi|\varphi}$}. Voici le résultat

a) $\unsure{M|\mathcal M}(\Phi)$ est la somme directe (à définir !) des $M(\Phi_i)$, où les $\Phi_i$ sont les composantes connexes de $\Phi$.

\deleted{$M(\Phi$} Si $\Phi$ \inserted{\unsure{est|ou}} connexe, \inserted{et non de la forme $\{a\}$, $a\in I$} $M(\Phi)$ est de la forme $M(\deleted{\unsure{I}}J,\Omega_0,\Omega_1)$, où $J,\Omega_0,\Omega_1$, sont définis ainsi, en termes de $\Phi$ :
% slanted
\margin{le cas $\Phi=\{a\}$ est trivial, $M\unsure{\Rightarrow|\simeq}$ \unsure{la maquette ponctuelle} (\unsure{zéro-dimensionnelle})}
\margin{$a=\operatorname{or}\delta\Phi$\\$b=\operatorname{ex}\delta\Phi$\\\deleted{$\unsure{\Phi=\Phi}$}}
\[
J=\Sigma_{\Phi}\cap\unsure{\mathscr M_0|\mathcal M_0|\mathbb M_0|M_0}
=\operatorname{Omb}^{0}_{I}(\Phi)\deleted{\unsure{a}}
=\left\{s\in I\ \middle|\
\begin{aligned}
&\text{a) }\deleted{\unsure{s}}a\leq s\leq b,\\
&\text{b) }s\vee\delta\Phi\in\unsure{\operatorname{Div}|\operatorname{Div}_{r}}I\\
&\text{c) }s\in\{a,b\}\Longrightarrow s\in\unsure{\Phi_0|\Phi^{0}}.
\end{aligned}\right\}
\tag*{(\unsure{a|a'})}
\]
\margin{\uline{Attention} $\longrightarrow$ cette description \unsure{OK} si $\Phi$ ne comporte pas de \unsure{diviseurs incompatibles} $\unsure{I_{\leq a}}$ et $\unsure{I_{\geq b}}$}

\scanpage{28}
\mspage{24}
% archive sheet number: 27

Donc, si à chaque strate \inserted{$X$} on associe son support interne $|X|^{0}\subset I$, avec la relation d'ordre induite par $I$ (NB $|X|^{0}$ peut être vide \inserted{si $X\in\unsure{M_1|\mathcal M_1}$} !), alors $J$ est la somme \deleted{ordonnée des} \inserted{lexicographique} des $|X|^{0}$ $(X\in\unsure{\Phi|\varphi})$, $X$ suivant \unsure{dans} l'ordre $\leq$ de $M$. D'autre part, \deleted{$\Omega_0$ est vide ou un \unsure{point}}
% slanted
\margin{cette description tjrs valable}
\[
\Omega_0=
\left\{\begin{aligned}
&\varnothing\text{ si le premier élément de }\Phi\text{ n'a pas d'origine}\\
&\quad\text{i.e. est un }I_{<a},\text{ ou s'il a}\\
&\deleted{\operatorname{or}(\Phi)\cap\Phi}\quad\text{une origine }a\deleted{\unsure{\Phi}}\text{ et }a\notin\Phi\\
&a=\operatorname{or}(\Phi)\text{ sinon}
\end{aligned}\right.
\]
\[
\deleted{\Omega_1=\operatorname{ex}(\Phi)\cap}
\]
$\Omega_1$ défini dualement

\uline{NB} Il faut définir l'origine d'une préfigure, quand elle existe (i.e. quand sa première strate \deleted{\unsure{est un}} n'est pas un $I_{<a}$), et dualement, l'extrémité.

\begin{center}\rule{2.8cm}{0.4pt}\end{center}

Cas $M$ \uline{divisible}. Dans le cas $(I,\omega_0,\omega_1)$ cela signifie que pour $a<b$ $(a,b\in\unsure{J|I})$ $\exists c$ avec $a<c<b$. Dans le cas des $(I,\omega_0)$ resp. $(I,\omega_1)$, il faut ajouter que $\forall a\in I$, existe $b\in I$ t.q. $b<a$ resp $b>a$. Enfin, dans le cas de l'intervalle $I$ biouvert, il faut \inserted{alors} ajouter encore : $I\neq\varnothing$. (\uline{NB} pour dire que l'\uline{ensemble ordonné} $I$ est divisible : donc $\varnothing$ ne sera pas considéré comme ens. ordonné divisible, bien que $M(\varnothing)=\{\varnothing\}$, magasin ponctuel, le soit !)

Si $I$ est divisible, alors l'ens $M_c\subset M$ des strates ``compactes'' (i.e. de la forme $a$ et $I_{ab}$) est très dense.

\scanpage{29}
\mspage{25}
% archive sheet number: 28

\uline{Digression} 1\quad Somme directe de magasins.

Si $(\unsure{M_i|\mathcal M_i})_{i\in I}$ est une famille de magasins, on définit le magasin somme $M$ ainsi
\[
M=\coprod_i M_i
\tag*{a)}
\]
b) $\trianglelefteq$ sur $M$ est la relation d'ordre somme

c) $\overset{\circ}{\ll}$ sur $M$ est la relation d'ordre somme

d) Si $x,y\in M$, on a
\[
x\mathrel{|\circ|}y\Longleftrightarrow
\left\{\begin{aligned}
&x\in M_i,y\in M_j,\ i\neq j\\
&\text{\uline{ou}}\\
&x,y\in M_i,\quad x\mathrel{|\circ|}y\\
&\text{au sens de }M_i
\end{aligned}\right.
\]

Vérification des axiomes tautologique.

Si pour chaque $i\in I$, on a choisi $\unsure{\mathcal F_i|\mathscr F_i|F_i}\in\unsure{P(M_i)|\mathcal P(M_i)}$, atelier associé à $M_i$, on définit
\[
\begin{aligned}
&\mathcal F\subset P(M)\\
&\rotatebox{90}{$=$}\\
&\{A\subset M\mid\forall i\in M_i,\quad A\cap M_i\in\mathcal F_i\}
\end{aligned}
\]
On voit que c'est un atelier associé à $M$, appelé \uline{atelier somme} des ateliers $(M_i,\trianglelefteq,\overset{\circ}{\ll},|\circ|,\mathcal F_i)$. On note que
\[
\mathcal F\simeq\prod_{i\in I}\mathcal F_i
\]
isomorphisme respectant les ordres $\unsure{\leq|\ll}$ $\subseteq,\ll$.

\scanpage{30}
\mspage{26}
% archive sheet number: 29

On fera attention que cette notion de somme n'est \uline{pas} une notion de somme au sens catégorique (pour la notion de (co)morphisme de magasins introduite jusqu'à présent. \unsure{Ce ne serait pas normal|Ce qui serait peu normal}, vu l'interprétation géométrique et le renversement des flèches -- \unsure{intuitivement|intentionnel|interstices}, la somme de magasins correspond bien à une \uline{somme} d'espaces (et non à un produit, comme pourrait le suggérer le principe de renversement des flèches).
% slanted
\margin{vu la condition de compatibilité avec $|\circ|$}
% slanted
\margin{\deleted{\unsure{semble}}}

Si on part d'un magasin $\unsure{M|\mathcal M}$, on peut songer à définir ses composantes connexes $M_i$, comme les composantes connexes pour la relation $\ll$ : On trouve une décomposition en somme, à condition de supposer l'axiome suivant (qui sera \unsure{toujours|beaucoup} \unsure{confirmé|employé} par la suite)

\margin{\rule{0.35pt}{1.5cm}}
\uline{Mag conn$^{\unsure{0|\circ}}$}\quad Si $X,Y\in M$ appartiennent à des comp. connexes distinctes de $M$, on a
\[
X\mathrel{|\circ|}Y.
\]

\scanpage{31}
\mspage{27}
% archive sheet number: 30

\uline{Digression 2}\quad Notion de sous-magasin.

Soit $\unsure{M|\mathcal M}$, \inserted{$\trianglelefteq,\overset{\circ}{\ll},|\circ|$} un magasin, $M'$ une partie, munissons la des relations induites -- quand peut-on dire que c'est un magasin.

Mag 1 trivial

Mag $2'$ a) aussi. Pour b, il faut supposer

[\uuline{SousMag}] Si $X\overset{\circ}{\ll}Y$ dans $M'$, et $X'\unsure{\triangleleft|\trianglelefteq}X$ dans $M'$, alors l'unique $Y'\in\widetilde Y$ tel que $X'\overset{\circ}{\ll}Y'$ est $\in M'$.
% slanted
\margin{\uline{NB} Il faudrait imposer une condition plus forte, à voir. \uline{Sousmag} Si $X\overset{\circ}{\ll}Y'\unsure{\triangleleft|\trianglelefteq}Y$, et $X,Y\in M'$, alors $Y'\in M'$, laquelle s'écrit (pour \inserted{$X,Y\in M'$}) : si $X\unsure{\overset{\circ}{\ll}|\ll}Y$ dans $M$, \deleted{alors} \inserted{alors $X\underset{M'}{\ll}Y$}. Les \unsure{couples plong. ds $M'$ sont alors compatibles|couples plong. ds $M'$ sont alors complets}. \unsure{serait plus forte}.}

(\uline{NB} c'est une figure si $M'$ fermé par $\trianglelefteq$)

Mag $3'$ trivial, Mag $4'$ trivial.

\uline{Ex 1} Soit $\unsure{\Phi|\varphi}$ \inserted{$\subset M$} une préfigure de $M$, i.e. une partie telle que $\overline\Phi^{(\trianglelefteq)}$ soit une figure, ou encore, \deleted{$\Phi$ fermé} $\forall X,Y\in\Phi$, on a $X\unsure{\gtrless|\lessgtr|\bowtie}Y$.

Considérons $\operatorname{Omb}^{0}(\Phi)$ (attention, c'est fermé pour $\overset{\circ}{\ll}$, mais pas en général pour $\trianglelefteq$). \deleted{\unsure{Dans ce cas}} \inserted{Alors} $M'=\operatorname{Omb}^{0}(\Phi)$ est un sous-magasin, i.e. satisfait Sousmag

\deleted{est un sous-magasin, i.e. satisfait la condition Sous-mag ci-dessus, il faut \unsure{s.} que l'on ait}
\[
\text{\deleted{\uline{Sousmag $\Phi$}}}\quad
\deleted{\forall Z\in\Phi\text{ et }X\overset{\circ}{\ll}Z,\text{ soit }X'\in\widetilde X,}
\]
\deleted{et $Z'$ son image dans $\widetilde Z_{\unsure{a|\alpha}}$ (donc $Z'\in\overline\Phi$), on veut que $X'\in\operatorname{Omb}^{0}(\Phi)\Rightarrow Z'\in\operatorname{Omb}^{0}(\Phi)$ i.e. en fait \uline{$Z'\in\Phi$ (puisque $Z'\in\overline\Phi$)} (même s'il n'est pas fermé par $\trianglelefteq$, quand $\Phi$ n'est pas une figure\ldots)}

\scanpage{32}
\mspage{28}
% archive sheet number: 31

En effet, soient $X,Y,X'\in\operatorname{Omb}^{0}(\unsure{\Phi|\varphi})$ avec
\diagram{\begin{tikzcd}[ampersand replacement=\&,column sep=0.7em,row sep=0.2em]
X\&\overset{\circ}{\ll}\&Y\\
\rotatebox{90}{$\trianglelefteq$}\&\&\\
X'\&\&
\end{tikzcd}}
alors $X'\ll Y\ll\overline\Phi$ donc $X'\ll\overline\Phi$, donc $\exists!\ Z'_0\in\overline\Phi$ avec $X'\overset{\circ}{\ll}Z'$. Comme $X'\in\operatorname{Omb}^{0}(\Phi)$, i.e. $\exists\ Z'_1\in\Phi$ tel que $X'\overset{\circ}{\ll}Z'_1$, on a $Z'_1=Z'_0$ i.e. $Z'_0\in\Phi$. Ceci posé, il existe puisque $Y\in\operatorname{Omb}^{0}(\Phi)$ un $Z\in\Phi$ avec $Y\overset{\circ}{\ll}Z$ d'où
\diagram{\begin{tikzcd}[ampersand replacement=\&,column sep=0.7em,row sep=0.2em]
X\&\overset{\circ}{\ll}\&Y\&\overset{\circ}{\ll}\&Z\in\Phi\\
\rotatebox{90}{$\trianglelefteq$}\&\&\rotatebox{90}{$\trianglelefteq$}\&\&\rotatebox{90}{$\trianglelefteq$}\\
X'\&\overset{\circ}{\ll}\&Y'\&\overset{\circ}{\ll}\&Z'\\
\&\&\rotatebox{90}{$\unsure{\in}$}?\&\&\\
\&\&\operatorname{Omb}^{0}(\Phi)\&\&
\end{tikzcd}}
; et on aura $Z'\in\overline\Phi$, donc $Z'=Z'_0\in\Phi$, donc $Y'\in\operatorname{Omb}^{0}(Z')\subset\operatorname{Omb}^{0}(\Phi)$, \unsure{qed.}

\uline{Ex 2} Soit $\Phi$ comme précédemment, et soit $\unsure{M'|\mathcal M'}=\unsure{\operatorname{Supp}|\operatorname{supp}}^{0}\Phi$ (un support constructible

Est-ce un sous-magasin ? C'est pratiquement équivalent à la $=$ question, avec $M'=\operatorname{cosupp}^{0}\Phi$. \deleted{(du moins, si} C'est faux en général, comme on voit sur la figure suivante.

\figure{\begin{tikzpicture}[scale=1.2,line cap=round,line join=round]
\coordinate (a) at (1.3,2.4);
\coordinate (b) at (0,0);
\coordinate (c) at (2,0);
\coordinate (d) at (1.2,0.08);
\draw[line width=1.1pt] (a)--(b);
\draw[line width=0.6pt] (a)--(c)--(b);
\draw[line width=0.8pt] (a)--(d);
\draw[line width=1.2pt] (b)--(d);
\begin{scope}
\clip (a)--(b)--(d)--cycle;
\foreach \k in {0.1,0.35,0.6,0.85,1.1,1.35,1.6,1.85}{\draw[line width=0.4pt] (0,\k)--(1.35,{\k+0.9});}
\end{scope}
\fill (d) circle (0.13);
\fill (a) circle (0.035);
\node[above] at (a) {$a$};
\node[left] at (b) {$b$};
\node[right] at (c) {$c$};
\node[below] at (d) {$d$};
\node[anchor=north west] at (3,2.6) {$\begin{aligned}
\deleted{\unsure{X}}Y&=2\text{-simplexe }(a,b,c)\\
X&=2\ \text{\rule{1.8cm}{0.35pt}}\ (a,b,d)\\
X'&=1\text{-simplexe }(b,d)\\
&\deleted{\Phi=\{d\}}\\
Y'&=1\ \text{\rule{1.8cm}{0.35pt}}\ (b,c)
\end{aligned}$};
\end{tikzpicture}}
$\Phi=\{d\}$. On a $\Phi\mathrel{|\circ|}Y$, \inserted{$\Phi\mathrel{|\circ|}X$,} (et $\Phi\mathrel{|\circ|}X'$)

\scanpage{33}
\mspage{29}
% archive sheet number: 32

\margin{\rule{0.35pt}{1.8cm}\hspace{3pt}\rule{0.35pt}{1.8cm}\hspace{3pt}\rule{0.35pt}{1.8cm}}
mais on n'a pas $\unsure{\Phi|\varphi}\mathrel{|\circ|}Y'$. Donc, attention à la définition du magasin résiduel d'une figure \inserted{$F$} (ici $\{d\}$, figure ponctuelle !).

Il faut, si on se borne aux multistrates $X$ avec $X\parallel F$, \unsure{pour rester sous|pour être stable sous} $\unsure{\ll|\trianglelefteq}$ dans un sous-magasin, \uline{se borner aux $X$ tels que $X\unsure{\gtrless|\lessgtr|\bowtie}F$}, ou $X\gtrless\Phi$ dans le cas d'une préfigure. (NB dans le cas de la dimension $1$, exceptionnellement, $X\mathrel{|\circ|}\Phi\Longrightarrow X\gtrless\Phi$, car $X\mathrel{|\circ|}Y\Longrightarrow X\gtrless Y$ ....). Je dis que si on prend
\[
\unsure{M'|\mathcal M'}=\{X\in M\mid X\mathrel{|\circ|}\Phi,\ X\gtrless\Phi\}\qquad ?
\]
% slanted
\margin{\uline{magasin résiduel large de $\Phi$.}}
alors $M'$ est un sous-magasin, i.e.
\diagram{\begin{tikzcd}[ampersand replacement=\&,column sep=0.7em,row sep=0.2em]
X\&\overset{\circ}{\ll}\&Y\\
\rotatebox{90}{$\trianglelefteq$}\&\&\rotatebox{90}{$\trianglelefteq$}\\
X'\&\overset{\circ}{\ll}\&Y'
\end{tikzcd}}
alors $X,Y,X'\in M\Longrightarrow Y'\in M$.

Comme $Y\gtrless\Phi$ et $Y'\trianglelefteq Y$, on a $Y'\gtrless\Phi$.

Donc pour voir $Y'\mathrel{|\circ|}\Phi$, il suffit de voir $Y'\notin\Phi$. Mais si on avait $Y'\in\Phi$, on aurait $X'\mathrel{|\circ|}Y'$, et comme $X'\overset{\circ}{\ll}Y'$, $X'\unsure{\mathrel{|\circ|}}X'$, absurde.

